\section{Related Works}
\label{sec:related-work}

\paragraph{Multi-agent debate (MAD).}
%MAD has emerged as an effective inference-time paradigm for improving the reasoning capability of LLMs. 
Early work~\cite{du2024multiagentdebate:mad-2,chen2024reconcile:mad-1,liang2024encouraging:mad-3} shows that multiple agents iteratively critiquing and revising each other's answers can outperform single-agent inference on knowledge-intensive and reasoning tasks. 
Subsequently, \citet{liu2024groupdebate} partition agents into subgroups to reduce token cost, while \citet{lin2025enhancing,taubenfeld2025confidence,fu2025deep} incorporate self-reported confidence to weight or filter agents' contributions. More recently, \citet{zhu2026demystifying} identify viewpoint homogeneity as a primary cause of debate stagnation; \citet{zhang2026key} show that positionally dominant agents disproportionately determine the final outcome; \citet{tian2026multi} reveal that erroneous content accumulated across rounds can persistently mislead downstream agents; and \citet{nguyen2026hear} preserve diversity by selectively retaining historically dissenting messages. 
%These works %diagnose failure modes induced by \emph{who} speaks and \emph{what} is remembered, but 
%largely retain a fixed communication structure across rounds. 
%AR-MAD %treats the communication topology as a first-class control variable and 
%dynamically routes critiques, %based on the instantaneous system state, thereby 
%mitigating the induced information homogenization and error propagation. %by fixed interaction structures.

\paragraph{Communication topology.}
%A parallel line of research investigates how communication structure shapes multi-agent performance. In LLM-based systems, 
\citet{li2024improving} demonstrate that sparse topologies can match or exceed fully connected debate at substantially lower cost. %, motivating the use of bandwidth-limited graphs. 
Subsequently, %methods learn or design the topology itself: 
\citet{zhang2024g:g-designer} propose G-Designer, which uses graph neural networks to architect task-specific communication graphs; \citet{shen2025assemble} autoregressively generate agent topologies; and \citet{li2025adaptive} adaptively prune edges to balance accuracy and efficiency. Beyond LLM debate, \citet{hu2024learning} learn communication graphs end-to-end, \citet{sun2024t2mac} perform targeted message selection, \citet{guan2024efficient} aggregate information via self-supervision, and \citet{10720863} optimize communication at the semantic level. 
These approaches typically %optimize a %\emph{single} 
%topology, either offline or once per task, 
rely on training signals or learned controllers. In contrast, PEAR adapts the topology purely at inference.  %using the evolving debate state to route information 
%without any model fine-tuning or ground-truth supervision.